% !TeX root = closed-systems-from-comparison-completeness.tex
% ============================================================
\chapter{Connection-First Reconstruction of Geometry}
\label{ch:connection}
% ============================================================
\begin{chapterorientation}
\orientationitem{Input}{The transport stack, its quadratic curvature carrier, and the macroscopic compatibility boundary.}
\orientationitem{Role}{This chapter reconstructs geometry from intrinsic connection data rather than from an assumed background coordinate manifold.}
\orientationitem{Output}{The first-order connection class, smooth realization interface, and geometric handoff to phase curvature.}
\end{chapterorientation}
\section{Introduction}
Under the standing principle of closed-world admissibility
(\cref{sp:closed-world}), this chapter develops the connection-first
realization endpoint of the transport stack determined by
\cref{sp:transport}.
Recasting the transport-to-curvature results, it develops a
connection-first reconstruction that bridges the intrinsic stack to
\cref{ch:em}.
Conventional connection, bundle, and curvature notation used for the
realized comparison is standard \cite{KobayashiNomizu1963,Nakahara2003}.
\Cref{ch:bottom,ch:lift,ch:interface,ch:causality,ch:hilbert,ch:einstein}
determine the intrinsic geometric content of the
closed stack up to first-order gauge.
Closure determines quotient semantics.  Admissible state content is exactly
quotient content on
\[
	\Phys:=X/G.
\]
Thus no invariant state datum remains except what descends through the
orbit projection
\[
	\pi:X\to\Phys.
\]
The classification of enrichment then shows that every
non-quotient enrichment is exhausted by two loci, possibly with both
and with no third:
\begin{enumerate}[label=\textup{(\alph*)}]
	\item
	      object-level representative choice;
	\item
	      morphism-level transport.
\end{enumerate}
Object-level enrichment is noncanonical. Once quotient semantics has
been determined, the only possible intrinsic carrier of geometry is
therefore transport.
Spacetime interleaving next constructs the canonical two-parameter
object
\[
	ST:=\varprojlim_{k,t}\bigl(S_k/F_t^{(k)}\bigr),
\]
which is the first object in the stack on which spatial refinement and
temporal coarsening coexist in a single intrinsic inverse-limit
construction.
Finally, the obstruction theory proves that the first visible intrinsic
transport residue is the stabilized quadratic carrier
\[
	\mathcal K\simeq F^2/F^3,
\]
and that, under the inherited second-jet faithfulness condition,
nonzero stabilized square classes are equivalent to nonzero realized
curvature in any compatible smooth realization.
That condition is no longer free-standing: it is equivalent to the
ghost-free regime \textup{(T)}+\textup{(D)} and is discharged on the
canonical realization
(\cref{sec:interface-ghosts,sec:interface-detectability,sec:interface-canonical}).
Thus, in the ghost-free regime of the canonical realization, geometry
is no longer conjectural.  It is already determined there.  The
remaining question is to determine the form in which geometry is
intrinsic before auxiliary choices are made.
The answer is connection-first geometry.
Indeed, an atlas is first-order coordinate data.  But first-order
coordinate data cannot be intrinsic in a closed system: a coordinate
assignment requires representative choice in orbit fibres, hence lies in
the object locus, and the object locus has already been proved
noncanonical.  What survives canonically after quotient descent is not
coordinates but transport.  Accordingly, the intrinsic geometric object
determined by the stack is not a coordinate manifold but a connection-type
structure: a spacetime precursor, a family of first-variation arenas, a
first-order transport law defined only up to gauge, and a first
tensorial obstruction carried by the quadratic carrier
\[
	\mathcal K.
\]
The chapter proves, in order:
\begin{enumerate}[label=\textup{(\roman*)}]
	\item
	      that no coordinate atlas can be canonically reconstructed from the
	      closed comparison structure;
	\item
	      that the regular locus of \(ST\) carries a canonical family of
	      four-dimensional first-variation spaces and a first-order transport law;
	\item
	      that these data determine an intrinsic connection class on the
	      first-variation bundle;
	\item
	      that the first tensorial residue of this connection class is precisely
	      the quadratic carrier
	      \[
		      \mathcal K\simeq F^2/F^3;
	      \]
	\item
	      that the entire transport filtration
	      \[
		      F^1\supseteq F^2\supseteq F^3\supseteq\cdots
	      \]
	      is the intrinsic jet hierarchy of this connection-first geometry;
	\item
	      that whenever a compatible smooth realization exists, the intrinsic
	      connection-first geometry realizes as ordinary Lorentzian connection
	      geometry and, under the inherited second-jet faithfulness condition,
	      nonzero stabilized square classes in \(\mathcal K\) are equivalent
	      to nonzero realized curvature in that realization.
\end{enumerate}
Thus the conclusion is exact:
\medskip
\begin{center}
	\emph{coordinates belong to gauge, transport does not.}
\end{center}
\medskip
Geometry in a closed comparison system is intrinsically
connection-first.
% ============================================================
\section{The spacetime precursor and the first-variation family}
\label{sec:connection-precursor}
% ============================================================
Recall from Chapter~\ref{ch:causality} that the refinement tower and the
observer-induced temporal quotient filtration determine the spacetime
object
\[
	ST:=\varprojlim_{k,t}\bigl(S_k/F_t^{(k)}\bigr).
\]
By construction, \(ST\) is the first object in the stack in which the
two structural directions
\[
	k\uparrow\ \text{(spatial refinement)},
	\qquad
	t\uparrow\ \text{(temporal coarsening)}
\]
coexist inside one intrinsic inverse-limit object.
At this stage \(ST\) is canonical.  What remains observer-relative is
the observer reduction and the finite levels to which it descends.  No
coordinate structure has been introduced.
The causality chapter already isolates the locus on which first-order
geometric data are nondegenerate.
\begin{definition}[Regular locus]
	Define
	\[
		ST^{\reg}\subseteq ST
	\]
	to be the subset of points \(p\in ST\) for which the first-variation
	arena \(V_p\) exists and satisfies
	\[
		\dim V_p=4.
	\]
\end{definition}
\begin{proposition}[Pointwise first-variation geometry]
	\label{prop:connection-pointwise}
	For every \(p\in ST^{\reg}\), the stack canonically determines:
	\begin{enumerate}[label=\textup{(\roman*)}]
		\item
		      a real vector space \(V_p\) of dimension \(4\);
		\item
		      a symmetric bilinear form
		      \[
			      B_p:V_p\times V_p\to\mathbb R;
		      \]
		\item
		      under the causal-cone hypotheses of Chapter~\ref{ch:causality}, a
		      Lorentzian signature condition
		      \[
			      \sign(B_p)\in\{(1,3),(3,1)\}.
		      \]
	\end{enumerate}
\end{proposition}
\begin{proof}
	This is exactly the pointwise conclusion proved in
	Section~\ref{sec:lorentz} of Chapter~\ref{ch:causality}.  The
	definition of \(ST^{\reg}\) merely isolates the locus on which that
	conclusion holds.
\end{proof}
\begin{definition}[First-variation bundle]
	\label{def:connection-first-bundle}
	Define
	\[
		T^{(1)}ST^{\reg}
		:=
		\bigsqcup_{p\in ST^{\reg}}V_p
	\]
	with projection
	\[
		\pi_1:T^{(1)}ST^{\reg}\to ST^{\reg},
		\qquad
		\pi_1(v)=p
		\quad (v\in V_p).
	\]
\end{definition}
% ============================================================
\section{No intrinsic coordinates}
\label{sec:no-intrinsic-coordinates}
% ============================================================
The coordinate-manifold picture assumes canonical local coordinate
assignments determined purely by the intrinsic structure of the system.
We show that such assignments are structurally impossible in a closed
comparison system.
\begin{definition}[Intrinsic coordinate system]
	An \emph{intrinsic coordinate system} on a space \(Y\) constructed from
	the closed comparison structure is a rule assigning to every point
	\(p\in Y\) a neighborhood \(U_p\subseteq Y\) and a map
	\[
		\phi_p:U_p\to\mathbb R^n
	\]
	such that the assignment \(p\mapsto(U_p,\phi_p)\) is determined
	functorially from the closed comparison data and is invariant under all
	automorphisms of the comparison structure.
\end{definition}
\begin{theorem}[Universal obstruction to intrinsic coordinates]
	\label{thm:no-intrinsic-coordinates}
	Let
	\[
		(X,G,\pi:X\to\Phys=X/G)
	\]
	be the quotient system determined by closure.
	Then no intrinsic coordinate system exists on \(\Phys\) or on any
	geometric object constructed solely from quotient-level state content.
	Equivalently, every coordinate construction necessarily depends on
	noncanonical gauge data.
\end{theorem}
\begin{proof}
	Closure identifies admissible state content with quotient content
	descending through
	\[
		\pi:X\to\Phys.
	\]
	Hence any intrinsic coordinate system on a quotient-level geometric
	object must be determined functorially from the quotient datum
	\((X,G,\pi)\).
	Let \(\Aut(X,G,\pi)\) denote the automorphism group of this quotient
	datum.  Every element
	\[
		\alpha\in\Aut(X,G,\pi)
	\]
	induces an automorphism of \(\Phys\).  If an intrinsic coordinate system
	exists, then functoriality requires the coordinate assignment to be
	equivariant under every such \(\alpha\).  Thus for each \(p\in\Phys\)
	and each \(\alpha\in\Aut(X,G,\pi)\) the local coordinate maps satisfy
	\[
		\phi_{\alpha(p)}\circ \alpha = \phi_p
	\]
	on the common domain of definition.
	Now fix \(p\in\Phys\).  Because \(\pi:X\to\Phys\) is an orbit
	projection, the fibre \(\pi^{-1}(p)\) is a \(G\)-orbit.  The group \(G\)
	acts transitively on that fibre, and this transitively permutes all
	representative choices above \(p\).  The induced symmetries extend to
	local automorphisms of the quotient datum.  Therefore equivariance under
	\(\Aut(X,G,\pi)\) determines any canonically assigned first-order local
	coordinate data to be invariant under this transitive representative
	action.
	But invariance under transitive representative action eliminates all
	distinguished first-order labels attached to representatives.  In
	particular, no canonically chosen coordinate germ can separate local
	states by a representative-dependent first-order assignment.  An
	intrinsic coordinate system would therefore fail to distinguish nearby
	points by locally injective coordinate maps.
	This contradicts the defining property of a coordinate system.
	Hence no intrinsic coordinate system exists.
\end{proof}
\begin{corollary}[Coordinates are gauge]
	Any coordinate atlas on a geometry derived from the closed comparison
	system necessarily depends on auxiliary gauge choices.
\end{corollary}
\begin{proof}
	By \cref{thm:no-intrinsic-coordinates}, coordinates cannot be
	constructed functorially from the intrinsic data.  Any coordinate
	assignment therefore requires extra choices breaking the automorphism
	symmetry of the quotient system.
\end{proof}
% ============================================================
\section{Transport as first-order geometry}
\label{sec:transport-first-order}
% ============================================================
The morphism locus already carries transport before any obstruction
appears.  Chapters~\ref{ch:rotation} and \ref{ch:lift} show that a
representative lift is exactly representative choice together with a
flat \(G\)-connection on the protocol category.  The later bridge,
interface, and causality chapters show that the first stable failure of
this flat compositional regime appears in degree \(2\).
This is the formal pattern of a connection with curvature.  The present
section extracts that pattern intrinsically.
\begin{definition}[First-order transport law]
	A \emph{first-order transport law} on \(ST^{\reg}\) is an assignment,
	for nearby regular points \(p,q\in ST^{\reg}\), of linear maps
	\[
		\tau_{q\to p}:V_q\to V_p
	\]
	obtained from the morphism-locus transport by passage to first
	variation.
\end{definition}
\begin{proposition}[Transport is the unique intrinsic first-order residue]
	\label{prop:connection-transport-unique}
	On \(ST^{\reg}\), every intrinsic geometric datum not already exhausted
	by quotient semantics must arise from morphism transport and not from
	object-level structure.
\end{proposition}
\begin{proof}
	Closure determines quotient semantics, so quotient-level state
	content is already complete.
	The two-locus classification shows that every non-quotient
	enrichment is exhausted by representative choice and morphism-level
	transport, possibly with both and with no third locus.
	By \cref{thm:no-intrinsic-coordinates}, representative choice cannot
	yield intrinsic first-order geometry.
	Therefore any intrinsic first-order residue must come from the
	morphism-transport locus.
\end{proof}
\begin{proposition}[Functoriality to first order]
	\label{prop:connection-functoriality}
	The first-order transport law satisfies:
	\begin{enumerate}[label=\textup{(\roman*)}]
		\item
		      \[
			      \tau_{p\to p}=\id_{V_p};
		      \]
		\item
		      whenever the relevant comparisons are defined,
		      \[
			      \tau_{r\to p}
			      =
			      \tau_{q\to p}\circ\tau_{r\to q}
		      \]
		      to first order;
		\item
		      the first visible failure of strict transport composition is invisible
		      in degree \(1\) and appears first in the quadratic carrier
		      \[
			      \mathcal K\simeq F^2/F^3.
		      \]
	\end{enumerate}
\end{proposition}
\begin{proof}
	Items \textup{(i)} and \textup{(ii)} are the first-order shadow of the
	flat transport formalism of Chapters~\ref{ch:rotation} and
	\ref{ch:lift}: before obstruction appears, representative transport is
	flat and composes functorially.
	Item \textup{(iii)} is exactly the content of the later obstruction
	theory.  The first stable intrinsic failure of flat transport
	composition does not survive in degree \(1\); it appears first in the
	degree--\(2\) quotient
	\[
		F^2/F^3,
	\]
	as proved in the bridge, interface, and causality chapters.
\end{proof}
\begin{corollary}[Connection pattern]
	\label{cor:connection-pattern}
	The intrinsic first-order geometry of \(ST^{\reg}\) has the formal shape
	of a connection: first-order transport composes functorially, and its
	first tensorial obstruction occurs in degree \(2\).
\end{corollary}
\begin{proof}
	This is the combined content of
	\cref{prop:connection-transport-unique,prop:connection-functoriality}.
\end{proof}
% ============================================================
\section{Classification of intrinsic first-order geometry}
\label{sec:first-order-classification}
% ============================================================
\Cref{cor:connection-pattern} identifies morphism transport as the only
possible source of intrinsic first-order geometry.  We now sharpen this
to a classification theorem.
\begin{definition}[First-order geometric residue]
	A \emph{first-order geometric residue} on \(ST^{\reg}\) is any
	assignment of linear comparison data between nearby first-variation
	spaces
	\[
		V_p
	\]
	which is functorial under refinement and compatible with quotient
	semantics.
\end{definition}
\begin{theorem}[Classification of first-order geometric residues]
	\label{thm:first-order-classification}
	Every intrinsic first-order geometric residue on \(ST^{\reg}\) arises
	uniquely from morphism-locus transport.
	Equivalently, the space of intrinsic first-order geometric structures on
	\(ST^{\reg}\) is canonically identified with the space of transport laws
	induced by representative lifts.
\end{theorem}
\begin{proof}
	Let \(R\) be a first-order geometric residue.  By definition, \(R\)
	assigns linear comparison maps between nearby first-variation spaces
	\[
		V_p.
	\]
	Closure implies that quotient semantics already exhausts invariant
	state content.
	Therefore any additional structure is exhausted by representative
	choice and morphism-level transport, possibly with both and with no
	third enrichment locus.
	Object-level enrichment depends on representative choice and is
	therefore noncanonical.
	Since \(R\) is intrinsic by assumption, it cannot arise from that
	locus.
	Hence \(R\) must arise from morphism transport.
	Conversely, the representative-lift formalism of
	Chapter~\ref{ch:lift} constructs exactly such transport laws.  Passing
	to first variation produces linear comparison maps
	\[
		\tau_{q\to p}:V_q\to V_p
	\]
	between nearby fibres, and these satisfy the required functoriality and
	quotient-compatibility conditions.
	Thus
	\[
		R\longleftrightarrow \tau
	\]
	is a bijective correspondence.  Hence every intrinsic first-order
	geometric residue arises uniquely from morphism transport.
\end{proof}
\begin{corollary}[Uniqueness of intrinsic first-order geometry]
	Up to gauge, the only intrinsic first-order geometric structure on
	\(ST^{\reg}\) is the transport law induced by representative lifts.
\end{corollary}
\begin{proof}
	Immediate from \cref{thm:first-order-classification}.
\end{proof}
% ============================================================
\section{Local gauges and the intrinsic connection class}
\label{sec:intrinsic-connection-class}
% ============================================================
Since coordinates are gauge, first-order transport must be handled by
local gauge-fixing rather than by a canonical atlas.
\begin{definition}[Local flattening gauge]
	A \emph{local flattening gauge} at a point \(p\in ST^{\reg}\) is a local
	choice of representative linearization of the first-variation family
	near \(p\) in which the first-order transport law is written explicitly
	as an affine connection-type law on \(V_p\).
\end{definition}
\begin{proposition}[Existence and uniqueness up to gauge]
	\label{prop:connection-flattening-gauges}
	Every regular point \(p\in ST^{\reg}\) admits local flattening gauges,
	and any two such gauges differ by first-order gauge transformation.
\end{proposition}
\begin{proof}
	Existence is the local form of the representative-lift formalism.  The
	morphism locus supplies local transport data; choosing local
	representatives identifies these data with first-order comparison maps
	between the first-variation spaces.  This is exactly a local flattening
	gauge.
	If two such gauges are chosen, they differ by local change of
	representatives.  Chapter~\ref{ch:lift} proves that changing
	representatives acts on transport by the gauge modification law.
	Passing to first variation, the two first-order transport descriptions
	therefore differ by first-order gauge transformation.
\end{proof}
\begin{definition}[Intrinsic connection class]
	The \emph{intrinsic connection class} on \(ST^{\reg}\) is the gauge
	equivalence class of first-order transport laws on
	\[
		T^{(1)}ST^{\reg}
	\]
	induced by all local flattening gauges.
\end{definition}
\begin{theorem}[Connection-first reconstruction]
	\label{thm:connection-first-reconstruction}
	Under the standing principle \cref{sp:closed-world}, the closed-system
	stack canonically determines on \(ST^{\reg}\):
	\begin{enumerate}[label=\textup{(\roman*)}]
		\item
		      the first-variation bundle
		      \[
			      T^{(1)}ST^{\reg}
			      =
			      \bigsqcup_{p\in ST^{\reg}}V_p;
		      \]
		\item
		      the pointwise causal type and, under the causal-cone hypotheses, the
		      pointwise Lorentzian first-variation pairings \(B_p\);
		\item
		      an intrinsic connection class on \(T^{(1)}ST^{\reg}\), defined modulo
		      first-order gauge by comparison transport.
	\end{enumerate}
\end{theorem}
\begin{proof}
	Item \textup{(i)} is \cref{def:connection-first-bundle}.  Item
	\textup{(ii)} is \cref{prop:connection-pointwise}.
	For item \textup{(iii)}, \cref{thm:first-order-classification} shows
	that every intrinsic first-order geometric structure on \(ST^{\reg}\)
	arises uniquely from morphism transport.  Passing to first variation
	produces linear comparison maps
	\[
		\tau_{q\to p}:V_q\to V_p
	\]
	between nearby fibres.  By \cref{prop:connection-functoriality}, these
	maps compose functorially to first order.
	Local representative choices produce gauges in which the transport law
	is written explicitly as an affine connection.  By
	\cref{prop:connection-flattening-gauges}, different gauges differ by
	first-order gauge transformation.  Therefore the gauge-equivalence class
	of such transport laws is canonically determined.
	This is exactly the intrinsic connection class.
\end{proof}
% ============================================================
\section{Quadratic obstruction as curvature carrier}
\label{sec:quadratic-obstruction-curvature}
% ============================================================
We now identify the first tensorial residue of the connection-first
geometry.
\begin{theorem}[Quadratic obstruction is the intrinsic curvature carrier]
	\label{thm:connection-curvature-carrier}
	For the intrinsic connection class reconstructed in
	\cref{thm:connection-first-reconstruction}, the first tensorial
	transport obstruction is the stabilized quadratic carrier
	\[
		\mathcal K\simeq F^2/F^3.
	\]
	Equivalently, \(\mathcal K\) is the intrinsic curvature carrier of the
	connection-first geometry.
\end{theorem}
\begin{proof}
	By the bridge, interface, and causality chapters, the first stable
	failure of flat transport composition appears in the degree--\(2\)
	quotient
	\[
		F^2/F^3.
	\]
	No tensorial degree--\(1\) obstruction survives.
	The intrinsic connection class reconstructed in
	\cref{thm:connection-first-reconstruction} is, by definition, the
	first-order geometry carried by morphism-locus transport on the
	first-variation bundle.  Therefore its first tensorial obstruction is
	exactly the first stable tensorial failure of transport composition.
	But that failure is precisely the stabilized quadratic carrier
	\(\mathcal K\).
	Hence \(\mathcal K\) is the intrinsic curvature carrier of the
	connection-first geometry.
\end{proof}
\begin{corollary}[Quadratic determination of curvature]
	The intrinsic geometric hierarchy on \(ST^{\reg}\) begins as
	\[
		\text{first-order transport class}
		\;\Longrightarrow\;
		\mathcal K\simeq F^2/F^3,
	\]
	and, under the inherited second-jet faithfulness condition in any
	compatible smooth realization, nonzero stabilized square classes in
	\(\mathcal K\) are equivalent to nonzero realized curvature.
\end{corollary}
\begin{proof}
	The first step is \cref{thm:connection-curvature-carrier}.  For the
	realized branch, \cref{thm:interface-curvature} gives the inherited
	second-jet interface criterion: in any compatible smooth realization,
	nonzero stabilized square classes in \(\mathcal K\) are equivalent to
	nonzero realized curvature.
\end{proof}
% ============================================================
\section{The transport filtration as an intrinsic jet tower}
\label{sec:jet-tower}
% ============================================================
\Cref{sec:quadratic-obstruction-curvature,thm:connection-curvature-carrier}
identify the stabilized quadratic carrier
\[
	\mathcal K \simeq F^2/F^3
\]
as the first tensorial obstruction to first-order transport and hence as
the intrinsic curvature carrier of the connection-first geometry.
The natural next question is whether the higher filtration quotients
\[
	F^m/F^{m+1}
	\qquad (m\ge 1)
\]
also admit an intrinsic geometric interpretation.
The answer is affirmative: the full transport filtration is the
intrinsic jet hierarchy of the connection-first geometry.
\begin{definition}[Transport jet of order \(m\)]
	Let \(p\in ST^{\reg}\), and choose a local flattening gauge near \(p\).
	The \emph{transport jet of order \(m\)} at \(p\) is the order--\(m\)
	equivalence class of the local transport law under identification of
	two gauges when their transport laws agree through order \(m\).
\end{definition}
\begin{definition}[Intrinsic jet carrier]
	For each \(m\ge 1\), define the \emph{intrinsic jet carrier of order
		\(m\)} to be the graded transport quotient
	\[
		\gr^m(F):=F^m/F^{m+1}.
	\]
\end{definition}
\begin{lemma}[Order filtration lemma]
	\label{lem:order-filtration}
	Let two local flattening gauges induce transport laws \(T\) and \(T'\)
	which agree through order \(m-1\).  Then their first difference defines
	a class in
	\[
		F^m/F^{m+1}.
	\]
	Moreover, this class vanishes if and only if \(T\) and \(T'\) agree
	through order \(m\).
\end{lemma}
\begin{proof}
	Agreement through order \(m-1\) means exactly that all lower-order
	transport residues have been identified.  By definition of the
	filtration, the first nonvanishing difference then lies in \(F^m\).  Two
	transport laws agree through order \(m\) if and only if this difference
	lies in \(F^{m+1}\).  Therefore the first difference is represented
	canonically by a class in
	\[
		F^m/F^{m+1},
	\]
	and that class vanishes precisely when the two laws agree through order
	\(m\).
\end{proof}
\begin{lemma}[Gauge-independence of the order--\(m\) class]
	\label{lem:gauge-order-m}
	The class in
	\[
		F^m/F^{m+1}
	\]
	defined in \cref{lem:order-filtration} is invariant under change of
	local flattening gauge modulo lower-order jet data.
\end{lemma}
\begin{proof}
	A change of local flattening gauge alters the transport law by a gauge
	reparametrization.  The lower-order terms of this reparametrization are
	already absorbed by the identification of transport laws modulo the jet
	data through order \(m-1\).  The first surviving remainder is therefore
	well-defined modulo \(F^{m+1}\).  Hence the induced class in
	\[
		F^m/F^{m+1}
	\]
	depends only on the order--\(m\) transport defect and not on the chosen
	gauge representative.
\end{proof}
\begin{theorem}[Jet Tower Theorem]
	\label{thm:jet-tower}
	For every integer \(m\ge 1\), the graded piece
	\[
		\gr^m(F)=F^m/F^{m+1}
	\]
	is the intrinsic carrier of the order--\(m\) transport jet.
	More precisely:
	\begin{enumerate}[label=\textup{(\roman*)}]
		\item
		      the first-order carrier
		      \[
			      \gr^1(F)=F^1/F^2
		      \]
		      determines the connection class of the first-order transport law;
		\item
		      the second-order carrier
		      \[
			      \gr^2(F)=F^2/F^3
		      \]
		      determines the first tensorial obstruction to flatness, namely the
		      curvature carrier;
		\item
		      for each \(m\ge 3\), the graded piece
		      \[
			      \gr^m(F)=F^m/F^{m+1}
		      \]
		      determines the intrinsic order--\(m\) Taylor defect of transport
		      composition, modulo all lower-order transport data.
	\end{enumerate}
	Equivalently, the filtration
	\[
		F^1 \supseteq F^2 \supseteq F^3 \supseteq \cdots
	\]
	is the intrinsic jet tower of the connection-first geometry.
\end{theorem}
\begin{proof}
	We argue by induction on \(m\).
	For \(m=1\), the first-order transport law is the linear comparison
	transport
	\[
		\tau_{q\to p}:V_q\to V_p
	\]
	constructed from the morphism locus.  By
	\cref{thm:first-order-classification,thm:connection-first-reconstruction},
	this is exactly the intrinsic first-order geometric residue of the
	closed stack.  Hence
	\[
		\gr^1(F)=F^1/F^2
	\]
	is the carrier of the first transport jet.
	For \(m=2\), the statement is precisely
	\cref{thm:connection-curvature-carrier}: the first failure of strict
	first-order transport closure is invisible in degree \(1\) and appears
	first in
	\[
		F^2/F^3.
	\]
	This degree--\(2\) carrier is the intrinsic curvature carrier.
	Assume now that the statement has been proved through order \(m-1\), so
	that the lower quotients
	\[
		\gr^1(F),\dots,\gr^{m-1}(F)
	\]
	have already been identified with the intrinsic transport jets through
	order \(m-1\).
	Choose a local flattening gauge.  In that gauge the local transport law
	admits an expansion whose coefficients through order \(m-1\) are already
	determined by the lower graded carriers by the induction hypothesis.  If
	two such transport laws agree through order \(m-1\), then by
	\cref{lem:order-filtration} their first difference determines a class in
	\[
		F^m/F^{m+1},
	\]
	and this class vanishes exactly when the laws agree through order \(m\).
	By \cref{lem:gauge-order-m}, this class is invariant under change of
	local flattening gauge modulo lower-order jet data.  Therefore the
	order--\(m\) defect is intrinsically carried by
	\[
		F^m/F^{m+1}.
	\]
	Hence \(\gr^m(F)\) is the intrinsic carrier of the order--\(m\)
	transport jet.  This completes the induction.
\end{proof}
\begin{corollary}[Intrinsic Taylor hierarchy]
	\label{cor:intrinsic-taylor-hierarchy}
	The local transport law is determined, up to gauge, by the full tower of
	graded carriers
	\[
		\bigl\{\gr^m(F)\bigr\}_{m\ge 1}.
	\]
	In particular, the connection-first geometry admits an intrinsic formal
	Taylor expansion whose \(m\)-th coefficient is carried by
	\[
		F^m/F^{m+1}.
	\]
\end{corollary}
\begin{proof}
	Immediate from \cref{thm:jet-tower}.
\end{proof}
\begin{corollary}[Full differential-geometric reconstruction under convergence]
	Assume, in addition, that the intrinsic formal transport expansion
	defined by the tower
	\[
		\{\gr^m(F)\}_{m\ge 1}
	\]
	is convergent in local flattening gauges.  Then the closed comparison
	stack reconstructs not merely the connection class and curvature
	carrier, but the full local differential-geometric transport law.
\end{corollary}
\begin{proof}
	By \cref{cor:intrinsic-taylor-hierarchy}, the tower
	\[
		\{\gr^m(F)\}_{m\ge 1}
	\]
	determines the formal transport expansion to all orders.  If that
	formal expansion converges, it determines the local transport law
	itself.  Hence the full local differential-geometric structure is
	reconstructed.
\end{proof}
% ============================================================
\section{Smooth realization}
\label{sec:connection-smooth-realization}
% ============================================================
\Cref{thm:interface-curvature,thm:lorentz-forced} already prove that
smooth geometry is recognized
whenever the intrinsic transport algebra admits smooth realization.  We
now express that fact in the language of connection-first
reconstruction.
\begin{theorem}[Smooth realization of the connection-first geometry]
	\label{thm:connection-smooth-realization}
	Under the standing principle \cref{sp:closed-world}, assume that the
	intrinsic transport geometry on \(ST^{\reg}\) admits a compatible
	smooth realization.  Then there exist:
	\begin{enumerate}[label=\textup{(\roman*)}]
		\item
		      a smooth \(4\)-dimensional manifold \(M\);
		\item
		      a Lorentz metric \(g\) on \(M\), with pointwise signature determined by
		      the realized first-variation pairings;
		\item
		      an affine connection \(\nabla\) on \(TM\);
	\end{enumerate}
	such that:
	\begin{enumerate}[label=\textup{(\alph*)}]
		\item
		      the intrinsic first-order transport class realizes as the connection
		      class of \(\nabla\);
		\item
		      under the inherited second-jet faithfulness condition, nonzero
		      stabilized square classes in \(\mathcal K\) are equivalent to
		      nonzero realized curvature in any compatible smooth realization;
		\item
		      under the metric realization already identified in the stack, this
		      curvature agrees with the Riemann curvature tensor of \(g\).
	\end{enumerate}
\end{theorem}
\begin{proof}
	By hypothesis, the intrinsic transport geometry admits a compatible
	smooth realization.  By \cref{thm:christoffel-main}, the intrinsic
	first-order transport class then realizes as the connection class of
	an affine connection on the realized regular locus.
	The pointwise first-variation data already determine the dimension and
	Lorentzian signature of the realized tangent model on the regular
	locus.  Thus the realization is \(4\)-dimensional and Lorentzian.
	By \cref{thm:connection-curvature-carrier}, the stabilized quadratic
	carrier \(\mathcal K\) is the intrinsic curvature carrier of the
	connection-first geometry.  The intrinsic--smooth interface statement in
	\cref{thm:interface-curvature} then yields only the inherited
	second-jet criterion: under that hypothesis, nonzero stabilized
	square classes in \(\mathcal K\) are equivalent to nonzero realized
	curvature in any compatible smooth realization.  This is exactly item
	\textup{(b)}.
	In the metric realization already identified in the stack, the
	realized curvature of that affine connection agrees with the Riemann
	curvature tensor of the realized Lorentz metric, giving item
	\textup{(c)}.
\end{proof}
% ============================================================
\section{Interpretation}
\label{sec:connection-interpretation}
% ============================================================
The chapter may be summarized in one sentence:
\medskip
\centerline{\emph{In a closed system, intrinsic geometry is
		connection-first.}}
\medskip
The exclusion and the construction are equally rigid.
An intrinsic atlas is impossible.  Coordinates require representative
choice and therefore belong to gauge.  The closed-system theorems do not
merely fail to reconstruct coordinates; they prove that coordinates are
not admissible intrinsic primitives.
What is determined is transport.  Once the representative sector has been
removed, the regular spacetime precursor carries a canonical
first-variation family, a canonical connection class, a canonical
curvature carrier, and indeed a full intrinsic jet hierarchy.
Thus the intrinsic geometric output of the stack is the chain
\[
	\begin{aligned}
		 & ST^{\reg}
		\;\Longrightarrow\;
		\{V_p\}_{p\in ST^{\reg}}
		\;\Longrightarrow\;
		\text{intrinsic connection class}
		\\
		 & \;\Longrightarrow\;
		\mathcal K\simeq F^2/F^3
		\;\Longrightarrow\;
		\{\gr^m(F)\}_{m\ge 1}.
	\end{aligned}
\]
A coordinate chart appears only after gauge-fixing.  The transport class
and its full filtration tower exist before that choice.
% ============================================================
\section{Structural synthesis}
\label{sec:connection-conclusion}
% ============================================================
The geometric endpoint of the closed relational stack is now exact.
The stack does not intrinsically reconstruct coordinates, and it does
not need to.  Coordinates are impossible as canonical objects in a
closed comparison system because they depend on representative choice
and therefore belong to the noncanonical object locus.
What the stack does reconstruct is the intrinsic transport geometry of
the regular spacetime precursor.
More precisely, it determines:
\begin{enumerate}[label=\textup{(\roman*)}]
	\item
	      the canonical spacetime object
	      \[
		      ST=\varprojlim_{k,t}\bigl(S_k/F_t^{(k)}\bigr);
	      \]
	\item
	      its regular locus \(ST^{\reg}\);
	\item
	      the pointwise first-variation arenas
	      \[
		      V_p
		      \qquad (p\in ST^{\reg}),
	      \]
	      each of dimension \(4\);
	\item
	      the pointwise Lorentzian causal type under the causal-cone hypotheses;
	\item
	      an intrinsic connection class on the first-variation bundle, defined
	      modulo first-order gauge by comparison transport;
	\item
	      the stabilized quadratic carrier
	      \[
		      \mathcal K\simeq F^2/F^3,
	      \]
	      which is the intrinsic curvature carrier of that geometry;
	\item
	      the full filtration tower
	      \[
		      F^1\supseteq F^2\supseteq F^3\supseteq\cdots,
	      \]
	      whose graded pieces
	      \[
		      F^m/F^{m+1}
	      \]
	      are the intrinsic jet carriers of the transport law.
\end{enumerate}
Accordingly, the true geometric spine of the theory is
\[
	\begin{aligned}
		(U,\mathcal C) & \Longrightarrow \text{closure}                   \\
		               & \Longrightarrow \text{quotient semantics}        \\
		               & \Longrightarrow \text{transport}                 \\
		               & \Longrightarrow ST                               \\
		               & \Longrightarrow \text{connection-first geometry} \\
		               & \Longrightarrow \mathcal K \simeq F^2/F^3        \\
		               & \Longrightarrow \{\text{gr}^m(F)\}_{m \ge 1}.
	\end{aligned}
\]
Whenever a compatible smooth realization exists, this intrinsic
transport geometry realizes as ordinary Lorentzian connection geometry,
and the full filtration realizes as the local differential-geometric
jet hierarchy.  Under the inherited second-jet faithfulness condition,
nonzero stabilized square classes in \(\mathcal K\) are equivalent to
nonzero realized curvature in any compatible smooth realization.
That is the correct endpoint of the reconstruction program at the
present stage of the manuscript.
\section{The local Einstein arena}
\label{sec:einstein-local-arena}
% ============================================================
\Cref{thm:connection-first-reconstruction,thm:jet-tower,thm:connection-smooth-realization}
reconstruct the intrinsic transport geometry of
the regular spacetime precursor
\[
	ST^{\reg}
\]
from the closed comparison stack.  The reconstruction proceeds without
the introduction of coordinates or background manifolds.  What is
obtained intrinsically is the following hierarchy:
\begin{enumerate}[label=\textup{(\roman*)}]
	\item
	      the spacetime precursor
	      \[
		      ST=\varprojlim_{k,t}\bigl(S_k/F_t^{(k)}\bigr);
	      \]
	\item
	      the regular locus \(ST^{\reg}\);
	\item
	      the first--variation arenas
	      \[
		      V_p \qquad (p\in ST^{\reg}),
	      \]
	      each a real vector space of dimension \(4\);
	\item
	      the intrinsic connection class induced by comparison transport;
	\item
	      the stabilized quadratic carrier
	      \[
		      \mathcal K \simeq F^2/F^3,
	      \]
	      which is the curvature carrier of this transport geometry.
\end{enumerate}
These objects constitute the full intrinsic geometric output of the
closed relational stack prior to any gauge choice.
\medskip
If the transport geometry admits a compatible smooth realization, the
connection-first structure reconstructed on \(ST^{\reg}\) realizes as
ordinary Lorentzian geometry.  More precisely, there exist
\begin{enumerate}[label=\textup{(\roman*)}]
	\item
	      a smooth four--dimensional manifold \(M\),
	\item
	      a Lorentz metric \(g\) on \(M\),
	\item
	      an affine connection \(\nabla\) on \(TM\),
\end{enumerate}
such that the intrinsic transport class realizes as the connection
class of \(\nabla\), and, under the inherited second-jet interface
hypothesis, nonzero stabilized square classes in the quadratic carrier
are equivalent to nonzero realized curvature for that connection.
In particular, the pointwise first--variation arenas \(V_p\) realize as
the tangent spaces \(T_pM\).  Since each \(V_p\) is four--dimensional,
the realized manifold is four--dimensional as well.
\medskip
Standard differential topology then implies that every point
\(p\in M\) admits a neighborhood \(U\subseteq M\) and a coordinate
chart
\[
	\phi:U\to\mathbb{R}^4
\]
which identifies \(U\) smoothly with an open subset of \(\mathbb{R}^4\).
Thus the final geometric output of the reconstruction is the familiar
local Einstein arena:
\[
	M \;\simeq\; \mathbb{R}^4
	\qquad \text{locally}.
\]
\medskip
The conceptual direction is therefore inverted relative to the usual
presentation of relativistic physics.
In the standard formulation of general relativity, one begins by
postulating a smooth four--dimensional manifold \(M\) equipped with a
Lorentz metric \(g\), and the theory proceeds from that geometric
background.
In the present framework, the manifold is not an assumption.  It is the
smooth realization of a deeper transport geometry reconstructed from
closed comparison structure.  Coordinates arise only after a gauge
choice, and the local model \(\mathbb{R}^4\) appears only at the final
stage of the reconstruction.
\medskip
The logical chain is therefore
\[
	\begin{aligned}
		 & (U,\mathcal C)
		\;\Longrightarrow\;
		\text{closure}
		\;\Longrightarrow\;
		\text{transport}
		\;\Longrightarrow\;
		ST^{\reg}
		\\
		 & \;\Longrightarrow\;
		\text{connection geometry}
		\;\Longrightarrow\;
		(M,g,\nabla),
	\end{aligned}
\]
with the classical Einstein starting point emerging as the final local
realization
\[
	M \;\simeq\; \mathbb{R}^4
	\quad \text{locally}.
\]
The Einstein manifold is thus not the beginning of the theory, but its
endpoint.

\section{Conclusion}
\label{sec:connection-transition}
The connection-first reconstruction is therefore complete at the stated
strength: the intrinsic transport geometry is reconstructed
unconditionally, while manifold and metric geometry appear as realized
endpoints only when a compatible smooth realization exists, not as
primitive background assumptions.
Accordingly, \cref{ch:em} extends this same intrinsic chain to the electromagnetic
phase sector.

\begin{chapterclosing}
\closingitem{Established}{Connection-first geometry is reconstructed from intrinsic transport; smooth manifold and metric geometry enter only as compatible realized endpoints.}
\closingitem{Carried forward}{The reconstructed connection-side structure supplies the setting for closed phase curvature.}
\end{chapterclosing}
